\chapter{Шестой спектральный момент циклической голономии и no-go масштаба}
\label{ch:v7-cycle-holonomy-spectral-moment-scale-gate}

\section{Оставшийся циклический объект}

После относительного семейного quotient полный ненулевой граф имеет один
независимый цикл
\begin{equation}
 C=(Q_L,u_R,X_L,e_R,L_L,Y_R,Q_L).
 \label{eq:v7-cycle-moment-cycle}
\end{equation}
Три старых ребра имеют фиксированную единичную семейную ориентацию, четыре
новых циклических ребра имеют общий радиальный масштаб $r$, а две
вектороподобные массы и древесное ребро $Q_Ld_R$ не входят в $C$.
После древесной фиксации всю относительную ориентацию можно поместить в один
блок $W_C\in U(3)$.

Пусть $D(r,W_C)$ является самосопряжённой блочной смежностью полного графа:
\begin{equation}
 D_{vw}=a_{vw}U_{vw},\qquad
 D_{wv}=a_{vw}U_{vw}^\dagger,\qquad
 a_{vw}=\begin{cases}1,&vw\in E_0,\\r,&vw\in E_*,\end{cases}
 \label{eq:v7-cycle-moment-dirac}
\end{equation}
где все древесные $U_{vw}=I_3$, а произведение вдоль $C$ равно $W_C$.

\section{Первый чувствительный спектральный момент}

Прямой подсчёт замкнутых маршрутов даёт
\begin{align}
 \Tr D^2&=18(2r^2+1),\nonumber\\
 \Tr D^4&=6(18r^4+10r^2+5),
 \label{eq:v7-cycle-moment-low-traces}
\end{align}
и обе формулы не зависят от $W_C$. Шестой момент равен
\begin{equation}
 \Tr D^6
 =54+144r^2+306r^4+324r^6
 +12r^4\Re\Tr W_C.
 \label{eq:v7-cycle-moment-sixth-trace}
\end{equation}
При $W_C=I_3$ это принимает форму
\begin{equation}
 \Tr D^6\big|_{W_C=I_3}
 =18(18r^6+19r^4+8r^2+3).
 \label{eq:v7-cycle-moment-sixth-trivial}
\end{equation}
Коэффициент двенадцать равен числу ориентированных начальных чтений простого
шестирёберного цикла: шесть выборов начальной вершины и две ориентации.
Следовательно,
\begin{equation}
 k_{\min}^{\rm hol}=6,\qquad
 \Delta\Tr D^6=12r^4(\Re\Tr W_C-3).
 \label{eq:v7-cycle-moment-first-sensitive-degree}
\end{equation}
Это согласуется с общим разложением спектрального действия колчана по
замкнутым путям \cite{v7-cycle-moment-perez}. След голономии является
классовой функцией и не зависит от выбора корневого семейного кадра.

\section{Что выбирает шестой момент}

Если коэффициент шестого момента в спектральном полиноме равен $c_6$, то
первая голономно-чувствительная часть есть
\begin{equation}
 V_6(W_C)=12c_6r^4\Re\Tr W_C.
 \label{eq:v7-cycle-moment-holonomy-potential}
\end{equation}
Её минимумы центральны:
\begin{equation}
 W_C^*=\begin{cases}
 I_3,&c_6<0,\\
 -I_3,&c_6>0,
 \end{cases}
 \qquad
 c_6=0:\quad W_C\ \text{не виден}.
 \label{eq:v7-cycle-moment-central-minima}
\end{equation}
Для $W_C=W_C^*\exp(iH)$ и ортонормированной следовой метрики на
$\mathfrak u(3)$ получаем
\begin{equation}
 \Hess V_6\big|_{W_C^*}=12|c_6|r^4I_9,
 \qquad (n_-,n_0,n_+)=(0,0,9).
 \label{eq:v7-cycle-moment-holonomy-hessian}
\end{equation}
Таким образом, ненулевой $c_6$ действительно поднимает девятимерный
линейный остаток предыдущего гейта. Но он выбирает только $\pm I_3$, а не
нецентральную матрицу смешивания. Знак ответа определяется спектральным
профилем, а не графом.

\section{Радиальный тест масштаба}

Возьмём общий усечённый полином
\begin{equation}
 \mathcal S_{246}(r,W_C)
 =c_2\Tr D^2+c_4\Tr D^4+c_6\Tr D^6.
 \label{eq:v7-cycle-moment-spectral-polynomial}
\end{equation}
После выбора центрального минимума при $c_6>0$ радиальная часть имеет вид
\begin{equation}
 \mathcal S_{246}(r)=\mathcal S_{246}(0)
 +ar^2+br^4+cr^6,
 \label{eq:v7-cycle-moment-radial-polynomial}
\end{equation}
где
\begin{equation}
 a=36c_2+60c_4+144c_6,\qquad
 b=108c_4+270c_6,\qquad
 c=324c_6.
 \label{eq:v7-cycle-moment-radial-coefficients}
\end{equation}
Ненулевой стационарный масштаб обязан удовлетворять
\begin{equation}
 a+2br^2+3cr^4=0.
 \label{eq:v7-cycle-moment-scale-equation}
\end{equation}
Фон $H_{15}$ вывел целые коэффициенты замкнутых маршрутов
$36,60,144,108,270,324$, но не числа $c_2,c_4,c_6$. Последние являются
моментами или производными выбранной спектральной функции и данными масштаба
усечения \cite{v7-cycle-moment-spectral-action}. Поэтому значение $r=\mu$
не следует из одного графа $H_{15}$.

Если все три коэффициента неотрицательны, то
\begin{equation}
 a>0,\qquad b>0,\qquad c\ge0,\qquad r_*=0,
 \label{eq:v7-cycle-moment-positive-profile-origin}
\end{equation}
то есть старый фон стабилизирует нуль, а не запускает конденсат. Ненулевой
масштаб требует отдельного знака спектрального профиля, уровня отображения
момента или нового динамического поля. В колчанных quotient постоянный сдвиг
уровня является параметром устойчивости/Fayet--Iliopoulos, если его не
выводит внешняя геометрия \cite{v7-cycle-moment-fi}.

\begin{theorem}[Спектральная видимость и центральное замыкание]
На текущем полном графе единственная относительная $U(3)$-голономия впервые
входит в спектральные моменты при степени шесть и имеет точный вклад
$12r^4\Re\Tr W_C$. При ненулевом коэффициенте шестого момента все девять
линейных голономных мод получают положительный гессиан, но минимум всегда
централен: $W_C=I_3$ или $W_C=-I_3$. Нецентральное семейное смешивание этим
моментом не выводится.
\end{theorem}

\begin{nogocriterion}[Фон не выводит уровень отображения момента]
Нельзя отождествлять комбинаторные коэффициенты замкнутых маршрутов с
коэффициентами спектрального профиля. Без независимого вывода
$c_2,c_4,c_6$ уравнение \eqref{eq:v7-cycle-moment-scale-equation} не
предсказывает $\mu$. Нельзя также выбирать знак $c_6$ по желаемой
голономии или называть центральный минимум выводом CKM/PMNS.
\end{nogocriterion}

\begin{protocol}
\begin{enumerate}
 \item независимых циклов полного графа: \textbf{$1$};
 \item первый чувствительный момент: \textbf{$\Tr D^6$};
 \item точный голономный вклад: \textbf{$12r^4\Re\Tr W_C$};
 \item $\Tr D^2$ и $\Tr D^4$ видят $W_C$: \textbf{нет};
 \item девять линейных мод при $c_6\ne0$: \textbf{подняты};
 \item минимум при $c_6<0$: \textbf{$I_3$};
 \item минимум при $c_6>0$: \textbf{$-I_3$};
 \item нецентральный минимум: \textbf{нет};
 \item CKM/PMNS: \textbf{не выведены};
 \item коэффициенты радиального полинома: \textbf{вычислены};
 \item $H_{15}$ определяет $c_2,c_4,c_6$: \textbf{нет};
 \item масштаб $\mu$ выведен: \textbf{нет};
 \item итог: \textbf{положительная спектральная видимость, центральный
 holonomy-only минимум и no-go масштаба}.
\end{enumerate}
\end{protocol}

Машинный сертификат сохранён в
\path{s2t_v7_cycle_holonomy_spectral_moment_scale_gate_results.json}.

\begin{thebibliography}{9}
\bibitem{v7-cycle-moment-perez}
C.~I.~P\'erez-S\'anchez,
\emph{The Spectral Action on Quivers}, arXiv:2401.03705.
\bibitem{v7-cycle-moment-spectral-action}
M.~Eckstein, B.~Iochum,
\emph{Spectral Action in Noncommutative Geometry},
SpringerBriefs in Mathematical Physics 27 (2018); arXiv:1902.05306.
\bibitem{v7-cycle-moment-fi}
D.~Robles-Llana,
\emph{Quivers, Quotients, and Duality}, arXiv:hep-th/0405230.
\end{thebibliography}